\documentclass[11pt,a4paper]{article}
\usepackage[T1]{fontenc}
\usepackage[utf8]{inputenc}
\usepackage{lmodern}
\usepackage{amsmath,amssymb,amsthm}
\usepackage[margin=1.15in]{geometry}
\usepackage[hidelinks]{hyperref}
\emergencystretch=3em

\newtheorem{theorem}{Theorem}[section]
\newtheorem{proposition}[theorem]{Proposition}
\newtheorem{lemma}[theorem]{Lemma}
\newtheorem{corollary}[theorem]{Corollary}
\theoremstyle{definition}
\newtheorem{remark}[theorem]{Remark}

\newcommand{\Z}{\mathbb{Z}}
\newcommand{\Q}{\mathbb{Q}}
\newcommand{\F}{\mathbb{F}}
\newcommand{\CC}{\mathbb{C}}
\newcommand{\Zp}{\Z_{(p)}}
\newcommand{\RP}{\mathbb{RP}}
\newcommand{\CP}{\mathbb{CP}}
\newcommand{\Kar}{\operatorname{Kar}}
\newcommand{\Hom}{\operatorname{Hom}}
\newcommand{\Ext}{\operatorname{Ext}}
\newcommand{\Tor}{\operatorname{Tor}}
\newcommand{\GL}{\operatorname{GL}}
\newcommand{\id}{\mathrm{id}}
\newcommand{\sk}[2]{#1^{#2}}

\title{Classifying spaces of finite monoids with bounded homology\thanks{AI-produced.
The main theorems were proved and the text was written by Claude Fable~5.2
(Anthropic) at effort max; the case of finite-dimensional classifying spaces and
the two examples are due to runs of GPT-6 Astra (OpenAI, ChatGPT Pro mode); see
Section~\ref{sec:provenance} for who proved, wrote and checked what. No human
has checked this text. To the extent that any rights exist in it, they are
dedicated to the public domain under CC0~1.0.}}
\author{}
\date{25 November 2026}

\begin{document}
\maketitle

\begin{abstract}
The classifying space $BM$ of a monoid $M$ is the geometric realization of the
nerve of $M$, regarded as a category with one object; its fundamental group is
the universal group $U(M)$, the group generated by the elements of $M$ with
the multiplication table of $M$ as relations. Let $M$ be finite. We prove that
if the integral homology of $BM$ vanishes in all large degrees, then $U(M)$
is trivial; more precisely, if a prime $p$ divides the order of $U(M)$, then
$H_k(BM;\F_p)\neq 0$ for infinitely many $k$. It follows that three conditions
on a finite monoid are equivalent: $BM$ is homotopy equivalent to a finite CW
complex; it is homotopy equivalent to a finite-dimensional one; its homology
is bounded. In that case $BM$ is simply connected. Together with a theorem of
Ben-Bassat, that every finite simply connected CW complex is homotopy
equivalent to $BM$ for a finite monoid $M$, this determines these homotopy
types: they are the finite simply connected complexes. A conjecture of the same
paper, that every connected finite CW complex with finite fundamental group is
of the form $BM$ with $M$ finite, is therefore false; the real projective plane
is a counterexample. The proof compares the $K$-theory of $BM$ with that of the
classifying space of a finite group contained in $M$, through Atiyah's
completion theorem. Two examples are worked out: a monoid with six elements
whose classifying space is simply connected and has homology in every degree
from $2$ on (it is a wedge of spheres), and a monoid with $23$ elements whose
classifying space is the complex projective plane.
\end{abstract}

\section{Introduction and results}\label{sec:intro}

\paragraph{Conventions.} A monoid $M$ is regarded as a category with one
object. Its nerve $NM$ is the simplicial set whose $k$-simplices are the
$k$-tuples $(m_1,\dots,m_k)$ of elements of $M$; the face $d_0$ drops $m_1$,
the face $d_i$ ($0<i<k$) replaces $m_i,m_{i+1}$ by the product $m_im_{i+1}$,
the face $d_k$ drops $m_k$, and the degeneracies insert $1$. The classifying
space $BM=|NM|$ is a CW complex with one $0$-cell and one $k$-cell for each
$k$-tuple of elements different from $1$; if $M$ is finite it has finitely many
cells in each dimension. We write $H_k(M;A)=H_k(BM;A)$; this group is also
$\Tor_k^{\Z M}(\Z,A)$ \cite[Proposition~2.2, Corollary~2.3]{Sw}. The homology of
$M$ is \emph{bounded} if $H_k(M;\Z)=0$ for all sufficiently large $k$. The
\emph{universal group} $U(M)$ is the group with one generator $\eta(m)$ for
each $m\in M$ and the relations $\eta(m)\eta(m')=\eta(mm')$; the map
$\eta\colon M\to U(M)$ is the universal homomorphism from $M$ to a group. A
space is simply connected if it is path connected with trivial fundamental
group. All CW complexes of type $K(G,1)$ of a group $G$ are denoted $BG$.

\begin{theorem}\label{thm:main}
Let $M$ be a finite monoid and $p$ a prime. If $H_k(M;\F_p)=0$ for all
sufficiently large $k$, then the order of $U(M)$ is prime to $p$. In
particular, if the homology of $M$ is bounded, then $U(M)$ is trivial.
\end{theorem}

\begin{theorem}\label{thm:equiv}
For a finite monoid $M$ the following conditions are equivalent:
\begin{itemize}
\item[\textup{(i)}] $BM$ is homotopy equivalent to a finite CW complex;
\item[\textup{(ii)}] $BM$ is homotopy equivalent to a finite-dimensional CW
complex;
\item[\textup{(iii)}] the homology of $M$ is bounded.
\end{itemize}
If they hold, $BM$ is simply connected.
\end{theorem}

For a finite group $G$, Theorem~\ref{thm:main} is contained in a theorem of
Swan: if $p$ divides $|G|$, then $H^i(G;\Z)$ has $p$-torsion for infinitely
many $i$ \cite[Corollary~1, and Corollary~2 for homology]{Swan}. For monoids
the fundamental group $U(M)$ can be trivial while $BM$ is any finite simply
connected complex. Fiedorowicz gave a monoid with five elements whose
classifying space is equivalent to the sphere $S^2$ (we return to it in
Section~\ref{sec:examples}) and suggested that every finite simply connected
complex should be equivalent to the classifying space of a finite monoid
\cite[p.~323]{Fie}. Ben-Bassat stated this as Conjecture~1.1 of \cite{BB}
and proved it:

\begin{theorem}[Ben-Bassat {\cite[Theorem~C]{BB}}]\label{thm:BB}
Every finite simply connected CW complex is homotopy equivalent to $BM$ for
some finite monoid $M$, which can be chosen aperiodic \textup{(}every $m\in M$
satisfies $m^{r+1}=m^r$ for some $r$\textup{)}.
\end{theorem}

\begin{corollary}\label{cor:class}
\textup{(a)} A connected finite CW complex is homotopy equivalent to $BM$ for
some finite monoid $M$ if and only if it is simply connected.

\textup{(b)} Up to homotopy equivalence, the spaces $BM$ for finite monoids
$M$ with bounded homology are exactly the finite simply connected CW
complexes.
\end{corollary}

\begin{corollary}\label{cor:conj}
Conjecture~7.1 of \cite{BB}, ``Let $X$ be a connected finite CW complex. If
$\pi_1(X)$ is finite, then $X\simeq BM$ for some finite monoid $M$'', is
false. The real projective plane $\RP^2$ is a counterexample, and so is every
connected finite CW complex whose fundamental group is finite and non-trivial.
\end{corollary}

If $BM$ is homotopy equivalent to a finite-dimensional complex, the
triviality of $U(M)$ can be deduced from a theorem of Miller on maps from
classifying spaces of torsion groups (Remark~\ref{rem:miller}); this was the
first route to Corollary~\ref{cor:conj}. Theorem~\ref{thm:main} replaces the
dimension by a condition on homology, which is what makes the equivalences of
Theorem~\ref{thm:equiv} possible.

\paragraph{What was known.} For a completely simple semigroup with an
identity adjoined, Steinberg computed the homology, and the homotopy type of
the classifying space, in terms of the structure group $G$ and the sandwich
matrix \cite[Theorem~A and Theorem~3.1]{St}; in that case bounded homology
forces $G=1$ by Swan's theorem, which is more than Theorem~\ref{thm:main}
says. Martinian and Steinberg proved that for a finite aperiodic regular
monoid the trivial module has a finite resolution by finitely generated
projective modules \cite[Corollary~5.15]{MS} (see also their Theorem~C and
Corollary~5.16); by Theorem~\ref{thm:equiv} the classifying space of such a
monoid, for instance of a finite band with an identity adjoined, is homotopy
equivalent to a finite simply connected complex. When the trivial module has
a projective resolution of finite length, the triviality of $U(M)$ has an
elementary proof by the restriction functor of Gray and Pride \cite{GP}, and
more is true (Corollary~\ref{cor:cd}). We have not found
Theorems~\ref{thm:main}, \ref{thm:equiv} and~\ref{thm:K}, or the failure of
Conjecture~7.1, in the literature; the searches are listed in the records
(Section~\ref{sec:provenance}).

Section~\ref{sec:K} proves a statement about maps from classifying spaces of
finite groups (Theorem~\ref{thm:K}); Section~\ref{sec:monoids} applies it to
monoids and proves Theorems~\ref{thm:main} and~\ref{thm:equiv};
Section~\ref{sec:class} contains the proofs of the two corollaries;
Section~\ref{sec:examples} gives the examples. The triviality of $U(M)$ does
not bound the homology (Proposition~\ref{prop:M6}), and rational homology
does not detect $U(M)$ (Remark~\ref{rem:rational}).

\section{Maps from classifying spaces of finite groups}\label{sec:K}

For a finite CW complex $Y$ let $K(Y)$ be the ring of virtual complex vector
bundles on $Y$, let $\sk{Y}{q}$ be the $q$-skeleton, and put
\[
 K_q(Y)=\ker\bigl(K(Y)\to K(\sk{Y}{q-1})\bigr)\quad(q\ge1),\qquad K_0(Y)=K(Y).
\]
We use two properties of this filtration. Both are contained in the spectral
sequence of Atiyah and Hirzebruch \cite[\S2]{AH}, which starts from
$H^*(Y;K^*(\mathrm{point}))$, converges to the graded group of the filtration
and is multiplicative; they are also in \cite[\S2]{At}, where the filtration
is defined and is multiplicative, and where (2.2) and Proposition~(2.4) give
its graded quotients.
\begin{itemize}
\item[(K1)] $K_a(Y)\cdot K_b(Y)\subset K_{a+b}(Y)$ for all $a,b\ge0$;
\item[(K2)] for $q\ge1$ the group $K_q(Y)/K_{q+1}(Y)$ is zero if $q$ is odd,
and is isomorphic to a subquotient of $H^q(Y;\Z)$ if $q$ is even.
\end{itemize}

Let $C$ be a cyclic group of order $p^a$, acting freely on the spheres
$S^{2j+1}\subset\CC^{j+1}$ through a faithful character, and let
$\sk{L}{2j+1}=S^{2j+1}/C$ be the lens spaces. Their union $L=S^\infty/C$ is a
$K(C,1)$ and has a CW structure with one cell in each dimension whose
$(2j+1)$-skeleton is $\sk{L}{2j+1}$; moreover $H^q(\sk{L}{2j+1};\Z)$ is
$\Z/p^a$ for $q$ even with $2\le q\le2j$ and is zero for $q$ odd with
$q<2j+1$, as one reads off from the cellular chain complex of
\cite[Example~2.43]{Ha}, whose boundary maps are alternately $0$ and
multiplication by $p^a$. Let $R(C)$ be the complex representation
ring. A representation $\tau$ of $C$ defines the vector bundle
$S^{2j+1}\times_C\tau$ on $\sk{L}{2j+1}$, and this gives ring homomorphisms
$\alpha_j\colon R(C)\to K(\sk{L}{2j+1})$ compatible with restriction.
\begin{itemize}
\item[(K3)] The homomorphism
$(\alpha_j)_j\colon R(C)\to\varprojlim_jK(\sk{L}{2j+1})$ is injective.
\end{itemize}
Indeed, by \cite[Proposition~(8.1)]{At} the completion of $R(C)$ at its
augmentation ideal maps isomorphically to the inverse limit of the $K$-groups
of the skeleta of a classifying space of the cyclic group $C$, and by
\cite[Proposition~(6.11)]{At} the map from $R(C)$ to its completion is
injective because $C$ is a $p$-group.

\begin{theorem}\label{thm:K}
Let $p$ be a prime, and let $X$ be a connected CW complex with finitely many
cells in each dimension such that $H^q(X;\Zp)=0$ for all $q>n$, where $\Zp$
is the ring of rational numbers with denominator prime to $p$. Let $G$ be a
finite group, $f\colon BG\to X$ a map, and
$\rho\colon\pi_1(X)\to\GL_d(\CC)$ a homomorphism. Then $\rho(f_*(g))=1$ for
every element $g\in G$ whose order is a power of $p$.
\end{theorem}

Here $f_*\colon G\to\pi_1(X)$ is defined up to conjugation in $\pi_1(X)$,
which does not affect the statement.

\begin{proof}
Let $g\in G$ have order $p^a$, let $C=\langle g\rangle$, and let $L$ and
$\sk{L}{2j+1}$ be as above. Since $BG$ is a $K(G,1)$, there is a map
$L\to BG$ inducing the inclusion $C\subset G$ on fundamental groups
\cite[Proposition~1B.9]{Ha}; composing with $f$ and applying cellular
approximation we get a cellular map $h\colon L\to X$ with $h_*=f_*|_C$ up to
conjugation. Fix base points, let $\widetilde X\to X$ be the universal
covering, on which $\pi_1(X)$ acts by deck transformations, and let
$V=\widetilde X\times_{\pi_1(X)}\CC^d$ be the vector bundle of rank $d$
defined by $\rho$. A lift $S^\infty\to\widetilde X$ of $h$ is equivariant
along $h_*$, so
\begin{equation}\label{eq:pullback}
 h^*V\cong S^\infty\times_C\sigma,\qquad\sigma=\rho\circ h_*\colon C\to\GL_d(\CC).
\end{equation}

\emph{Step 1.} For $m\ge0$ let $v_m\in K(\sk{X}{m})$ be the class of
$V|\sk{X}{m}$ minus the trivial bundle of rank $d$, and let $w_m=v_m^{n+1}$.
The skeleta are finite complexes, and restriction
$K(\sk{X}{m+1})\to K(\sk{X}{m})$ is a ring homomorphism taking $v_{m+1}$ to
$v_m$ and $w_{m+1}$ to $w_m$. The class $v_m$ has rank zero at every vertex,
so $v_m\in K_1(\sk{X}{m})$, and $w_m\in K_{n+1}(\sk{X}{m})$ by (K1). For
$m\le n$ this says $w_m=0$, since $K_{n+1}(\sk{X}{m})$ is then the kernel of
the identity of $K(\sk{X}{m})$.

\emph{Step 2: $w_m\otimes1=0$ in $K(\sk{X}{m})\otimes\Zp$ for every $m$.} We
argue by induction on $m$; the cases $m\le n$ are Step~1. Let $m>n$, assume
$w_{m-1}\otimes1=0$, and put $Y=\sk{X}{m+1}$, so that
$\sk{Y}{m-1}=\sk{X}{m-1}$ and $\sk{Y}{m}=\sk{X}{m}$. Since $\Zp$ is flat
over $\Z$, the kernel of $K(Y)\otimes\Zp\to K(\sk{Y}{q-1})\otimes\Zp$ is
$K_q(Y)\otimes\Zp$ for every $q$. The element $w_{m+1}\otimes1$ of
$K(Y)\otimes\Zp$ restricts to $w_{m-1}\otimes1=0$ on $\sk{Y}{m-1}$, hence
lies in $K_m(Y)\otimes\Zp$. By (K2) the group
$\bigl(K_m(Y)/K_{m+1}(Y)\bigr)\otimes\Zp$ is zero or a subquotient of
$H^m(Y;\Z)\otimes\Zp$. The cellular cochain complexes of $Y$ and of $X$
coincide in degrees at most $m+1$, so $H^m(Y;\Z)=H^m(X;\Z)$; and
$H^m(X;\Z)\otimes\Zp=H^m(X;\Zp)$, because the cellular cochain groups of $X$
are finitely generated and $\Zp$ is flat. This group is zero since $m>n$.
Therefore $K_m(Y)\otimes\Zp=K_{m+1}(Y)\otimes\Zp$, and $w_{m+1}\otimes1$
restricts to zero on $\sk{Y}{m}=\sk{X}{m}$; that is, $w_m\otimes1=0$.

\emph{Step 3.} The map $h$ sends $\sk{L}{2j+1}$ into $\sk{X}{2j+1}$. By
\eqref{eq:pullback} the pull-back of $v_{2j+1}$ is $\alpha_j(\sigma)-d$,
hence the pull-back of $w_{2j+1}$ is $\alpha_j\bigl((\sigma-d)^{n+1}\bigr)$,
and by Step~2
\[
 \alpha_j\bigl((\sigma-d)^{n+1}\bigr)\otimes1=0\quad\text{in }
 K(\sk{L}{2j+1})\otimes\Zp\quad\text{for every }j\ge0.
\]
The element $\alpha_j((\sigma-d)^{n+1})$ has rank zero, so it lies in
$K_1(\sk{L}{2j+1})$. By (K2) and the cohomology of lens spaces recalled
above, $K_1(\sk{L}{2j+1})$ has a finite filtration, ending with
$K_{2j+2}(\sk{L}{2j+1})=0$, whose quotients are zero or subquotients of
$\Z/p^a$; so it is a finite abelian $p$-group $A$, and $A\to A\otimes\Zp$ is
an isomorphism. Hence $\alpha_j((\sigma-d)^{n+1})=0$ for every $j$, and
$(\sigma-d)^{n+1}=0$ in $R(C)$ by (K3).

\emph{Step 4.} Taking characters, $(\chi_\sigma(g)-d)^{n+1}=0$, so
$\chi_\sigma(g)=d$. The matrix $\sigma(g)$ has finite order; its $d$
eigenvalues are roots of unity with sum $d$, so they are all equal to $1$,
and $\sigma(g)=1$. As $\sigma(g)$ is conjugate to $\rho(f_*(g))$, the theorem
follows.
\end{proof}

\begin{corollary}\label{cor:finitepi}
Let $X$ be a connected CW complex with finitely many cells in each dimension
and with finite fundamental group, let $G$ be a finite group and
$f\colon BG\to X$ a map.

\textup{(a)} If $H^q(X;\Zp)=0$ for all sufficiently large $q$, then the order
of $f_*(G)$ is prime to $p$.

\textup{(b)} If $H^q(X;\Z)=0$ for all sufficiently large $q$, then $f_*$ is
trivial.
\end{corollary}

\begin{proof}
(a) Apply Theorem~\ref{thm:K} to the regular representation $\rho$ of
$\pi_1(X)$, which is faithful: $f_*(g)=1$ for every $g$ of $p$-power order.
The image of a Sylow $p$-subgroup of $G$ is a Sylow $p$-subgroup of $f_*(G)$,
so the latter has order prime to $p$. (b) Since the cochain groups are
finitely generated, $H^q(X;\Zp)=H^q(X;\Z)\otimes\Zp$, so (a) applies to every
prime.
\end{proof}

\begin{remark}\label{rem:miller}
(a) Miller's theorem \cite[Theorem~B, p.~40; proved as Theorem~10.1,
pp.~84--85]{Mi} states that every map from $BG$, with $G$ a torsion group, to
a finite-dimensional CW complex is trivial on fundamental groups. Its proof
also uses the $K$-theory of cyclic groups, on a covering space of the target;
there is no hypothesis on the fundamental group of the target.
Theorem~\ref{thm:K} weakens the hypothesis on the dimension to one on
cohomology, and the argument of Step~2, on the skeleta of $X$ itself, takes
the place of the covering space. Ordinary cohomology cannot replace
$K$-theory here: under the surjection $\Z/p^2\to\Z/p$ a generator $y$ of
$H^2(B(\Z/p);\Z)$ goes to $p$ times a generator $x$ of $H^2(B(\Z/p^2);\Z)$,
so $y^k$ goes to $p^kx^k=0$ for $k\ge2$, and the map on integral cohomology
vanishes above degree~$2$, although the map is not trivial on fundamental
groups.

(b) Under the hypothesis on cohomology alone the representation $\rho$
cannot be dispensed with. By \cite[Theorem~A and the properties (1), (4)
recalled in the introduction]{PW} the group $\Z/2$ embeds in a group
$\Gamma$ which has the integral homology of a point and is of type $F_n$ for
every $n$, that is, has for every $n$ a classifying space with finite
$n$-skeleton. Such a group admits a classifying space $B\Gamma$ with finitely
many cells in each dimension. It is built skeleton by skeleton, starting
from a finite $2$-skeleton: if $Z$ is a finite complex of dimension $n\ge2$
with fundamental group $\Gamma$ whose universal covering $\widetilde Z$ is
$(n-1)$-connected, then the $\Z\Gamma$-module
$\pi_n(Z)\cong H_n(\widetilde Z;\Z)$ is finitely generated, by Schanuel's
lemma applied to the cellular chains of $\widetilde Z$ and of the universal
covering of a classifying space with finite $(n+1)$-skeleton, so that
finitely many $(n+1)$-cells attached to $Z$ make the universal covering
$n$-connected. Then $X=B\Gamma$ satisfies $H^q(X;\Z)=0$ for $q>0$, and $B(\Z/2)\to X$ is injective
on fundamental groups.
\end{remark}

\section{Finite monoids}\label{sec:monoids}

\begin{lemma}\label{lem:known}
Let $M$ be a finite monoid.

\textup{(a)} $\pi_1(BM)\cong U(M)$, the loop formed by the $1$-cell $(m)$
corresponding to $\eta(m)$.

\textup{(b)} $\eta\colon M\to U(M)$ is surjective and $U(M)$ is finite.

\textup{(c)} Let $e$ be an idempotent of the minimal two-sided ideal of $M$.
Then $H=eMe$ is a finite group with identity $e$, and $\eta$ restricts to a
surjective homomorphism of groups $H\to U(M)$.
\end{lemma}

\begin{proof}
(a) The nerve has a single vertex, so the fundamental group of $BM$ is
generated by the loops of the $1$-cells $(m)$, with one relation for each
$2$-cell $(m,m')$, which identifies the loop of $mm'$ with the composite of
the loops of $m$ and $m'$; this is the presentation of $U(M)$. The statement
is due to McDuff \cite[Lemma~1, p.~314]{McD}; it is also the case of one
object of
\cite[\S1, Proposition~1]{Qu}, by which the fundamental group of the
classifying space of a small category at an object is the automorphism group
of that object in the groupoid obtained by inverting all arrows. (b) is
proved in \cite[\S7]{BB}. (c) is known: the minimal ideal is a completely
simple semigroup, $H$ is its maximal subgroup at $e$, and $U(M)$ is the
quotient of $H$ by the normal closure of the entries of the sandwich matrix;
see \cite[\S2.1 and Lemma~2.1]{Sw}, with the attributions given there, and
\cite[p.~7]{St}. (Directly: for $x\in H$ the ideal
$MxM$ is contained in the minimal ideal, hence equal to it, so $e=axb$ with
$a,b\in M$ and $e=(eae)x(ebe)$; in a finite monoid one-sided inverses are
two-sided, so $x(ebe)(eae)=e$ and $x$ is invertible in $H$.) Since
$\eta(e)^2=\eta(e)$ we have $\eta(e)=1$; so $\eta|_H$ is a homomorphism of
groups, and $\eta(eme)=\eta(m)$ for $m\in M$ shows with (b) that it is
surjective.
\end{proof}

\begin{lemma}\label{lem:je}
With $e$ and $H$ as in Lemma~\ref{lem:known}\,\textup{(c)}, there is a map
$j_e\colon BH\to BM$ whose effect on fundamental groups is
$\eta|_H\colon H\to U(M)$, up to an inner automorphism.
\end{lemma}

\begin{proof}
We use the map of \cite[Introduction]{BB}. The Karoubi envelope $\Kar(M)$
is the category whose objects are the idempotents of $M$, with
$\Hom(f,g)=gMf$ and composition given by the product of $M$. The
endomorphism monoid of the object $1$ is $M$, that of $e$ is $H$; let
$\iota_1\colon M\to\Kar(M)$ and $\iota_e\colon H\to\Kar(M)$ be the
inclusions. Every object $f$ is a retract of $1$, by the arrows $1\to f$
and $f\to1$ given by the element $f$, so $B\iota_1$ is a homotopy
equivalence \cite[Lemma~2.4]{BB}. Let $j_e$ be $B\iota_e$ followed by a
homotopy inverse of $B\iota_1$.

Let $\Pi$ be the groupoid obtained from $\Kar(M)$ by inverting all arrows,
and write $[x]$ for the class of an arrow $x$; by \cite[\S1,
Proposition~1]{Qu} the fundamental groups of $B\Kar(M)$ at the objects $1$
and $e$ are the automorphism groups of $1$ and $e$ in $\Pi$. Let
$r\colon1\to e$ and $s\colon e\to1$ be the arrows given by the element $e$.
Then $r\circ s=ee=\id_e$, so $[s]=[r]^{-1}$. For $h\in H$,
\[
 [r]^{-1}[h][r]=[s\circ h\circ r]=[ehe]=[h],
\]
where on the right $h$ is regarded as an arrow $1\to1$, that is, as an
element of $M$. So the change of base point along $[r]$ carries the image of
$h\in H=\pi_1(BH)$ under $B\iota_e$ to the image under $B\iota_1$ of the loop
of $h\in M$, which is $\eta(h)$ by Lemma~\ref{lem:known}\,(a).
\end{proof}

\begin{proof}[Proof of Theorem~\ref{thm:main}]
Write $H_k=H_k(BM;\Z)$; these groups are finitely generated. Suppose that
$H_k(M;\F_p)=0$ for $k\ge k_0$. By the universal coefficient theorem
$H_k\otimes\F_p=0$ for $k\ge k_0$, so these $H_k$ are finite of order prime
to $p$. For $q>k_0$ the group
$H^q(BM;\Z)\cong\Hom(H_q,\Z)\oplus\Ext(H_{q-1},\Z)$ is then finite of order
prime to $p$, and $H^q(BM;\Zp)=H^q(BM;\Z)\otimes\Zp=0$. Thus $X=BM$ satisfies
the hypotheses of Corollary~\ref{cor:finitepi}\,(a): it is connected with
finitely many cells in each dimension, and its fundamental group $U(M)$ is
finite by Lemma~\ref{lem:known}. Applied to the map $j_e\colon BH\to BM$ of
Lemma~\ref{lem:je}, the corollary says that the image of $(j_e)_*=\eta|_H$
has order prime to $p$. This image is $U(M)$ by Lemma~\ref{lem:known}\,(c).

If the homology of $M$ is bounded, then $H_k(M;\F_p)\cong
H_k\otimes\F_p\oplus\Tor(H_{k-1},\F_p)$ vanishes for all large $k$ and every
prime $p$, so no prime divides the order of $U(M)$.
\end{proof}

\begin{proof}[Proof of Theorem~\ref{thm:equiv}]
(i) implies (ii), and (ii) implies (iii) because homology vanishes above the
dimension. Assume (iii). By Theorem~\ref{thm:main} and
Lemma~\ref{lem:known}\,(a), $BM$ is simply connected. Its homology groups
are finitely generated and vanish above some degree $n$. By
\cite[Proposition~4C.1]{Ha}, a simply connected CW complex is homotopy
equivalent to a CW complex having, for chosen decompositions of its homology
groups into cyclic summands, one $k$-cell for each infinite cyclic summand of
$H_k$ and one $k$-cell and one $(k+1)$-cell for each finite cyclic summand of
$H_k$. For $BM$ this complex is finite. This proves (i) and the last
assertion.
\end{proof}

\begin{corollary}\label{cor:swan}
Let $M$ be a finite monoid. If $U(M)\neq1$, then $H_k(M;\Z)\ne0$ for
infinitely many $k$. If a prime $p$ divides the order of $U(M)$, then
$H_k(M;\F_p)\neq0$ for infinitely many $k$.
\end{corollary}

For instance the monoid with $11$ elements of \cite[Example~3.4]{BB}, in the
case $a=1$ of that example, whose classifying space is $B(\Z/2)\vee S^2$,
has universal group $\Z/2$ and homology in infinitely many degrees.
Corollary~\ref{cor:swan} says that torsion coming from $U(M)$ recurs in
infinitely many degrees, as for groups. Torsion in higher degrees need not
recur: there are finite semigroups whose classifying spaces are Moore spaces
\cite[Table~1 and Examples~5.4, 5.5]{Sw}. The converse of
Corollary~\ref{cor:swan} fails. In the other case of the same example, where
$a$ is the generator $g$ of $\Z/2$, the classifying space is the cofibre of
the inclusion $\RP^1\subset\RP^\infty$ by \cite[Theorem~A]{BB}; it is simply
connected, and its reduced homology is $\Z$ in degree $2$, $\Z/2$ in every
odd degree from $3$ on, and zero in the other degrees. Another example is
the monoid $M_6$ of Proposition~\ref{prop:M6}.

\begin{corollary}\label{cor:cd}
Let $M$ be a finite monoid such that the trivial left \textup{(}or
right\textup{)} $\Z M$-module $\Z$ has a projective resolution of finite
length. Then the group $H=eMe$ of Lemma~\ref{lem:known}\,\textup{(c)} is
trivial. In particular $U(M)$ is trivial, and $BM$ is homotopy equivalent to
a finite simply connected CW complex.
\end{corollary}

\begin{proof}
We treat left modules; for right modules exchange the sides. Let $e$ and $H$
be as in Lemma~\ref{lem:known}\,(c). For a left $\Z M$-module $A$ the
subgroup $eA$ is a left $\Z H$-module, and $A\mapsto eA$ is a functor which
is exact, because $e$ is idempotent, commutes with direct sums and takes the
trivial module $\Z$ to the trivial module $\Z$; this is the functor of
\cite[Lemmas~2--4 and Theorem~2]{GP}. It takes $\Z M$ to $\Z[eM]$, and this
is a free $\Z H$-module, because $H$ acts freely on the set $eM$ by left
multiplication. Indeed, let $x\in eM$. Then $x$ lies in the minimal ideal, so
$e=axb$ with $a,b\in M$, and $(eae)(xbe)=e$ because $ex=x$; hence $xbe$ is
the inverse $g$ of $eae$ in $H$, and $xy=e$ for $y=beg^{-1}$. So $hx=h'x$
with $h,h'\in H$ implies $h=h'$. Therefore the functor takes free modules to
free modules, projective modules to projective modules, and a projective
resolution of finite length of $\Z$ over $\Z M$ to one over $\Z H$. For a
non-trivial finite group $H$ the groups $H^i(H;\Z)=\Ext^i_{\Z H}(\Z,\Z)$ are
non-zero for infinitely many $i$ \cite[Corollary~1]{Swan}; so $H=1$, and
$U(M)=\eta(H)=1$ by Lemma~\ref{lem:known}\,(c). Finally
$H_k(M;\Z)=\Tor_k^{\Z M}(\Z,\Z)$ can be computed from a projective resolution
of either variable, so the homology of $M$ is bounded, and
Theorem~\ref{thm:equiv} applies.
\end{proof}

\begin{remark}\label{rem:cd}
In Corollary~\ref{cor:cd} the triviality of $U(M)$ also follows from
Theorem~\ref{thm:main}, since the homology is bounded; and from Miller's
theorem (Remark~\ref{rem:miller}), since a monoid of finite cohomological
dimension has a classifying space of finite-dimensional homotopy type
\cite[Theorem~6.27 and Corollary~6.7]{GS}. We do not know whether $H=1$ for
every finite monoid with bounded homology. This holds for completely simple
semigroups with an identity adjoined \cite[Theorem~A]{St} and under the
hypothesis of Corollary~\ref{cor:cd}; Theorem~\ref{thm:main} only says that
$H$ has trivial image in $U(M)$; and the monoids of Theorem~\ref{thm:BB} can
be chosen aperiodic, so that $H=1$ for them. The map $j_e$ of
Lemma~\ref{lem:je} cannot decide the question: if the homology of $M$ is
bounded, then $BM$ is homotopy equivalent to a finite complex by
Theorem~\ref{thm:equiv}, and every map from $BH$ to a connected
finite-dimensional complex is null-homotopic, by Miller's theorem that the
space of pointed maps is weakly contractible \cite[Theorem~A, p.~39]{Mi}. An
answer has to use the monoid itself, as the proof of
Corollary~\ref{cor:cd} does.
\end{remark}

\section{The finite complexes that are classifying spaces of finite monoids}
\label{sec:class}

\begin{proof}[Proof of Corollary~\ref{cor:class}]
(a) If $X$ is a finite simply connected CW complex, $X\simeq BM$ with $M$
finite by Theorem~\ref{thm:BB}. Conversely, if a connected finite CW complex
$X$ is homotopy equivalent to $BM$ with $M$ finite, then $BM$ is simply
connected by Theorem~\ref{thm:equiv}, hence so is $X$. (b) By
Theorem~\ref{thm:equiv} every $BM$ with $M$ finite of bounded homology is
homotopy equivalent to a finite simply connected CW complex; conversely each
such complex is homotopy equivalent to some $BM$ with $M$ finite by
Theorem~\ref{thm:BB}, and the homology of this $M$ is bounded.
\end{proof}

\begin{proof}[Proof of Corollary~\ref{cor:conj}]
By Corollary~\ref{cor:class}\,(a) a connected finite CW complex with
non-trivial fundamental group is not homotopy equivalent to $BM$ for any
finite monoid $M$; the fundamental group of $\RP^2$ has order $2$.
\end{proof}

\begin{remark}\label{rem:onlyv1}
Corollary~\ref{cor:conj} concerns the first version of \cite{BB}, of 28
October 2026, the only one on arXiv on 25 November 2026; its Theorem~C is not
affected. Among finite categories with several objects the real projective
plane does occur: the category with objects $0,1,2$, with
$\Hom(i,j)=\Z/2$ for $i<j$, no other arrows than identities, and composition
given by addition, has nine arrows and a classifying space homeomorphic to
$\RP^2$, the quotient of the octahedron by the antipodal map
\cite[Theorem~6.2]{CC}; there this three-object model of $\RP^2$ is
attributed to Tanaka.
\end{remark}

\section{Examples}\label{sec:examples}

\begin{remark}\label{rem:rational}
Rational coefficients do not suffice in Theorem~\ref{thm:main}: for a finite
group $G\neq1$ one has $H_k(G;\Q)=0$ for all $k>0$ and $U(G)=G$.
\end{remark}

\subsection{A monoid with six elements: trivial universal group, unbounded
homology}

Let $M_6=\{1,x_{00},x_{01},x_{10},x_{11},y\}$. Put $r(x_{ij})=i$,
$c(x_{ij})=j$ and $r(y)=c(y)=0$; the element $1$ is the identity, and
\[
 vw=x_{r(v)\,c(w)}\qquad\text{for }v,w\neq1.
\]
Both ways of bracketing a product of three elements different from $1$ give
the row of the first and the column of the last, so $M_6$ is a monoid. It is
the semigroup of \cite[Example~5.2]{Sw} with an identity adjoined. The
elements $1$ and $x_{ij}$ form a submonoid $M_5$, which is the monoid with
five elements of Fiedorowicz \cite[p.~323]{Fie}; it is treated after the
proof of Proposition~\ref{prop:M6}. Write $e=x_{00}$.

\begin{proposition}\label{prop:M6}
\textup{(a)} $U(M_6)$ is trivial.

\textup{(b)} $H_0(M_6;\Z)=\Z$, $H_1(M_6;\Z)=0$ and $H_n(M_6;\Z)\cong\Z$ for
every $n\ge2$.

\textup{(c)} $BM_6$ is homotopy equivalent to the wedge
$\bigvee_{n\ge2}S^n$ of one sphere of each dimension $n\ge2$.
\end{proposition}

So the converse of Theorem~\ref{thm:main} fails, and $BM_6$ is a simply
connected space with finitely generated homology groups that satisfies none
of the conditions of Theorem~\ref{thm:equiv}. Part (b) is
\cite[Example~5.2]{Sw}.

\begin{proof}
(a) Every $x_{ij}$ is idempotent, so $\eta(x_{ij})=1$; and $ye=e$ gives
$\eta(y)=1$.

(b) Let $R=\Z M_6$ and
\[
 u=1-e,\quad a=x_{10}-e,\quad b=x_{01}-e,\quad
 c=x_{11}-x_{10}-x_{01}+e,\quad t=y-e .
\]
Then $e,u,a,b,c,t$ is a $\Z$-basis of $R$, the elements $e$ and $u$ are
idempotents with $eu=ue=0$, and the multiplication law gives
\[
 ae=a,\quad ea=0,\quad eb=b,\quad be=0,\quad ec=ce=et=te=0,\quad ab=c,
\]
while all other products of two of the elements $a,b,c,t$ are zero. Hence
$Re=\Z e\oplus\Z a$ and $Ru=\Z u\oplus\Z b\oplus\Z c\oplus\Z t$; these left
modules are projective. Consider
\[
 \cdots\xrightarrow{d_4}Re\oplus Ru\xrightarrow{d_3}Re\oplus Ru
 \xrightarrow{d_2}Ru\xrightarrow{d_1}Re\xrightarrow{\varepsilon}\Z\to0,
\]
with $\varepsilon$ the augmentation, $d_1(s)=sa$, $d_2(s,s')=sb+s't$ and
$d_n(s,s')=(0,sb+s't)$ for $n\ge3$. It is exact: the kernel of $\varepsilon$
on $Re$ is $\Z a=d_1(Ru)$, since $ua=a$; the map $d_1$ kills $b,c,t$, so its
kernel is $N=\Z b\oplus\Z c\oplus\Z t$; the map $s\mapsto sb$ sends $e,a$ to
$b,c$, and $s'\mapsto s't$ sends $u$ to $t$ and kills $b,c,t$; so $d_2$ has
image $N$ and kernel $0\oplus N$, and the same holds for $d_n$, $n\ge3$,
with image $0\oplus N$. Tensoring with the trivial right module, $Re$ gives
$\Z$ and $Ru$ gives $0$, because $\varepsilon(e)=1$ and $\varepsilon(u)=0$,
and all differentials become zero since the images of $d_n$, $n\ge2$, lie in
summands $Ru$. This gives (b).

(c) We use the cellular chain complex of $BM_6$: it is free on the tuples of
elements different from $1$, with
$\partial(m_1,\dots,m_n)=(m_2,\dots,m_n)
+\sum_{0<i<n}(-1)^i(\dots,m_im_{i+1},\dots)+(-1)^n(m_1,\dots,m_{n-1})$.
Fix $n\ge2$ and put
\[
 (a_i^0,a_i^1)=(x_{01},e)\ \ (i=1),\qquad(y,e)\ \ (1<i<n),\qquad
 (x_{10},e)\ \ (i=n).
\]
For $i<n$ the elements $a_i^\delta$ have row $0$, and for $i>1$ they have
column $0$; hence
\begin{equation}\label{eq:prod}
 a_i^{\delta}a_{i+1}^{\delta'}=e\qquad(0<i<n;\ \delta,\delta'\in\{0,1\}).
\end{equation}
Let $z_n=\sum_{\epsilon\in\{0,1\}^n}(-1)^{\epsilon_1+\dots+\epsilon_n}
(a_1^{\epsilon_1},\dots,a_n^{\epsilon_n})$. In $\partial z_n$ the terms
coming from the first, from the $i$-th and from the last summand of
$\partial$ do not depend on $\epsilon_1$, on $\epsilon_i$ (by
\eqref{eq:prod}) and on $\epsilon_n$ respectively, so they cancel in pairs:
$z_n$ is a cycle.

For $m\ne1$ let $\alpha(m)=1$ if $c(m)=1$ and $\alpha(m)=0$ otherwise, let
$\beta(m)=1$ if $r(m)=1$ and $\beta(m)=0$ otherwise, and let $Y(y)=1$ and
$Y(m)=0$ for $m\neq y$. For $m,m'\ne1$ we have $\alpha(mm')=\alpha(m')$,
$\beta(mm')=\beta(m)$ and $Y(mm')=0$. Define the cochain
\[
 f_n(m_1,\dots,m_n)=\alpha(m_1)\,Y(m_2)\cdots Y(m_{n-1})\,\beta(m_n).
\]
In $(\delta f_n)(m_0,\dots,m_n)$ the term $f_n(m_1,\dots,m_n)$ cancels
against $-f_n(m_0m_1,m_2,\dots)$, the last term against the term with
$m_{n-1}m_n$, and the remaining terms vanish because they contain a factor
$Y(m_im_{i+1})$. So $f_n$ is a cocycle, and $f_n(z_n)=1$, only
$\epsilon=0$ contributing. By (b) the class of $z_n$ is therefore a
generator of $H_n(BM_6;\Z)\cong\Z$.

This class is the image of the fundamental class of a sphere. Let $D_n$ be
the join of a circle with two vertices $p_0,p_1$ and two edges $E^0,E^1$ from
$p_0$ to $p_1$, and of $n-1$ two-point sets $\{q_i^0,q_i^1\}$, $2\le i\le
n$. Call $i$ the colour of $q_i^\delta$, and $0$ and $1$ the colours of $p_0$
and $p_1$. Then $D_n$ is a $\Delta$-complex whose simplices have vertices of
distinct colours, ordered by colour; its realization is the $(n-1)$-fold
suspension of a circle, a sphere $S^n$; its $n$-simplices are the joins
$\sigma_\epsilon$ of $E^{\epsilon_1}$ with $q_2^{\epsilon_2},\dots,
q_n^{\epsilon_n}$, and $\zeta_n=\sum_\epsilon(-1)^{\epsilon_1+\dots+
\epsilon_n}\sigma_\epsilon$ is a cycle, which generates $H_n(D_n)$ because
its coefficients are $\pm1$. Label the edges of $D_n$: the edge $E^\delta$
by $a_1^\delta$; an edge from a vertex of colour $i-1$ to $q_i^\delta$ by
$a_i^\delta$; every edge between vertices whose colours differ by at least
$2$ by $e$. For a triangle with vertices of colours $i<j<k$, the labels
$l,l'$ of the edges $i\to j$ and $j\to k$ have row $0$ and column $0$
respectively, as in \eqref{eq:prod}, so $ll'=e$, which is the label of the
edge $i\to k$. Hence sending a simplex with vertices
$v_0,\dots,v_k$ to the tuple of the labels of the edges
$v_0\to v_1,\dots,v_{k-1}\to v_k$ commutes with the face maps, and defines a
cellular map $s_n\colon S^n\cong|D_n|\to BM_6$ which sends the simplex
$\sigma_\epsilon$ onto the cell $(a_1^{\epsilon_1},\dots,a_n^{\epsilon_n})$;
so $(s_n)_*\zeta_n=z_n$.

All vertices go to the $0$-cell, so the maps $s_n$ define
$\bigvee_{n\ge2}S^n\to BM_6$. It induces an isomorphism on integral
homology, and both spaces are simply connected CW complexes, by (a) and
Lemma~\ref{lem:known}\,(a); so it is a homotopy equivalence
\cite[Corollary~4.33]{Ha}.
\end{proof}

\paragraph{The monoid with five elements.} In the submonoid
$M_5=\{1,x_{00},x_{01},x_{10},x_{11}\}$ of $M_6$ the product is
$x_{ij}x_{kl}=x_{il}$. This is the monoid of \cite[p.~323]{Fie}, where it is
left as an exercise that its classifying space is equivalent to $S^2$.
Fiedorowicz published a proof, by a projective resolution of length~$2$, in
\cite[Lemma, pp.~33--34]{Fie02}, and the case of every finite rectangular band
is in \cite[Theorem~3.1 and p.~8]{St}. The proof of
Proposition~\ref{prop:M6} contains another proof. The ring
$R'=\Z M_5$ has the $\Z$-basis $e,u,a,b,c$, with $R'e=\Z e\oplus\Z a$ and
$R'u=\Z u\oplus\Z b\oplus\Z c$, and the sequence
\[
 0\to R'e\xrightarrow{d_2}R'u\xrightarrow{d_1}R'e
 \xrightarrow{\varepsilon}\Z\to0,\qquad d_1(s)=sa,\quad d_2(s)=sb,
\]
is exact, because $d_2$ takes the basis $e,a$ of $R'e$ to the basis $b,c$ of
the kernel of $d_1$. Hence, as in the proof of (b), $H_2(M_5;\Z)\cong\Z$ and
$H_k(M_5;\Z)=0$ for $k\neq0,2$. The map $s_2$ takes $S^2$ into the subcomplex $BM_5$ of $BM_6$,
and $f_2(z_2)=1$; so $s_2\colon S^2\to BM_5$ induces an isomorphism on
homology and is a homotopy equivalence, as at the end of the proof above.
The element $y$ has the same products as $x_{00}$ with every element
different from $1$, without being equal to it; adjoining it to $M_5$ changes
the classifying space from $S^2$ to $\bigvee_{n\ge2}S^n$.

\subsection{A monoid with 23 elements of the homotopy type of the complex
projective plane}

Let $I=\{0,1,2\}$ and $J=\{0,1,2,3,4\}$. Consider pairs $(\sigma,\tau)$ of
maps $\sigma\colon I\to I$, $\tau\colon J\to J$ with the product
$(\sigma,\tau)(\sigma',\tau')=(\sigma\circ\sigma',\tau'\circ\tau)$. Let
$M_{23}$ consist of the identity, the $15$ pairs $x_{ij}$ of constant maps
with values $i$ and $j$, and the seven pairs below, where a map is written
as the string of its values:
\[
\begin{array}{c|ccccccc}
 &a&b&c&d&e&f&g\\\hline
\sigma&002&002&010&202&202&202&212\\
\tau&02200&41144&04034&03300&12211&40044&00334
\end{array}
\]
Write $\sigma_m,\tau_m$ for the components of $m\in M_{23}$. In this
subsection $e$ denotes the element of the table.

\begin{proposition}\label{prop:M23}
The set $M_{23}$ is closed under the product; it is an aperiodic monoid with
$23$ elements. The trivial left $\Z M_{23}$-module has a projective
resolution of length $4$, the groups $H_k(M_{23};\Z)$ are $\Z$ for
$k=0,2,4$ and zero otherwise, and $BM_{23}$ is homotopy equivalent to
$\CP^2$.
\end{proposition}

\begin{proof}
Closure is a finite verification. We have $x_{ij}x_{kl}=x_{il}$, the
elements $a,b,c,g$ are idempotent, and $d^2=x_{20}$, $e^2=x_{22}$,
$f^2=x_{24}$ with $d^3=d^2$, $e^3=e^2$, $f^3=f^2$; so $M_{23}$ is aperiodic.

Let $R=\Z M_{23}$. The elements $x_{00}$, $c-x_{00}$, $1-c$, $a-x_{00}$ of
$R$ are idempotents; let $P,E,F,G$ be the left ideals they generate, which
are projective. Define
\[
 0\to P\xrightarrow{d_4}G\oplus G\xrightarrow{d_3}P\oplus F
 \xrightarrow{d_2}E\xrightarrow{d_1}P\xrightarrow{\varepsilon}\Z\to0
\]
by the following values on the generating idempotents, extended
$R$-linearly (each value is fixed under left multiplication by the
idempotent it is assigned to, so the maps are well defined):
\[
\begin{aligned}
 d_1(c-x_{00})&=x_{10}-x_{00}, &
 d_2(x_{00},0)&=x_{03}-x_{00},\\
 d_2(0,1-c)&=c+d-g-x_{00}, &
 d_3(a-x_{00},0)&=(x_{00}-x_{20},\ b+e-f-x_{01}),\\
 d_3(0,a-x_{00})&=(0,\ a-x_{00}), &
 d_4(x_{00})&=(x_{02}-x_{00},\ x_{00}-x_{02}).
\end{aligned}
\]
The ranks over $\Z$ of $P,E,F,G$ are $3,10,10,4$, so the terms have ranks
$3,10,13,8,3$. All composites vanish, and the images of
$\varepsilon,d_1,d_2,d_3,d_4$ have ranks $1,2,8,5,3$ and are saturated
subgroups (the elementary divisors of the matrices are all $1$). Since
$3=1+2$, $10=2+8$, $13=8+5$, $8=5+3$ and $3=3$, each image has the rank of
the next kernel and, being saturated, equals it: the sequence is exact.
The augmentations of the four generating idempotents are $1,0,0,0$; so
tensoring with the trivial right module leaves $\Z$ in degrees $0,2,4$ and
$0$ in degrees $1,3$, which gives the homology. By Corollary~\ref{cor:cd},
$BM_{23}$ is homotopy equivalent to a finite simply connected complex, and
by \cite[Proposition~4C.1]{Ha} to $S^2\cup_he^4$ for some map
$h\colon S^3\to S^2$.

It remains to see that $h$ has Hopf invariant $\pm1$, that is, that the
square of a generator of $H^2$ generates $H^4$. Let
$\lambda\colon J\times I\to\Z$ be zero except for
$\lambda(1,1)=\lambda(1,2)=\lambda(2,2)=\lambda(4,1)=1$ and
$\lambda(3,1)=-1$, and put
\[
 u(m,m')=\lambda\bigl(\tau_m(0),\sigma_{m'}(0)\bigr)\qquad(m,m'\in M_{23}).
\]
Then $u$ vanishes if $m$ or $m'$ is $1$ and satisfies
$u(m',m'')-u(mm',m'')+u(m,m'm'')-u(m,m')=0$ for all $m,m',m''$; it is a
$2$-cocycle on $BM_{23}$, and $(m_1,\dots,m_4)\mapsto u(m_1,m_2)u(m_3,m_4)$
represents the square of its class. To evaluate these cocycles on generators
of $H_2$ and $H_4$ we map the resolution above to the normalized bar
resolution $B_\bullet$: here $B_n$ is the free left $R$-module on the symbols
$[m_1|\dots|m_n]$ with all $m_i\ne1$, the differential is
\begin{multline*}
 \partial(z[m_1|\dots|m_n])=zm_1[m_2|\dots|m_n]\\
 +\sum_{0<i<n}(-1)^iz[\dots|m_im_{i+1}|\dots]+(-1)^nz[m_1|\dots|m_{n-1}],
\end{multline*}
and $\Z\otimes_RB_\bullet$ is the cellular chain complex of $BM_{23}$. Let
$s$ be the additive map with $s(z[m_1|\dots|m_n])=[z|m_1|\dots|m_n]$ for
$z\in M_{23}$, $z\neq1$, and $s(1[m_1|\dots|m_n])=0$; then
$\partial s+s\partial=\id$ on $B_n$ for $n\ge1$. Put $F_0(x_{00})=x_{00}[\ ]$
and, for a generating idempotent $\epsilon$ of a summand of the term of
degree $n\ge1$, $F_n(\epsilon)=\epsilon\cdot s\bigl(F_{n-1}(d_n\epsilon)
\bigr)$, extended $R$-linearly. By induction $F$ is a chain map over the
identity of $\Z$, so it induces an isomorphism on homology after tensoring
with $\Z$. Let $z_2$ and $z_4$ be the images in $\Z\otimes_RB_\bullet$ of
$F_2(x_{00},0)$ and $F_4(x_{00})$; they are cycles representing generators
of $H_2$ and $H_4$. Evaluating gives
\[
 \langle u,z_2\rangle=-1,\qquad
 \langle u\smile u,z_4\rangle=1 .
\]
So the class of $u$ generates $H^2(BM_{23};\Z)$ and its square generates
$H^4$. Hence the Hopf invariant of $h$ is $\pm1$. Since
$\pi_3(S^2)$ is infinite cyclic, generated by the Hopf map
\cite[Example~4.51]{Ha}, whose Hopf invariant is $1$, and the Hopf invariant
is a homomorphism \cite[\S4.B, Proposition~4B.1]{Ha},
$h$ is homotopic to the Hopf map or to its negative, and
$S^2\cup_he^4\simeq\CP^2$.

The verifications called finite above (closure, the identities in $R$, the
ranks and elementary divisors, the cocycle identity, the chain map property
of $F$ and the two evaluations) were done by a program in exact integer
arithmetic; see Section~\ref{sec:provenance}.
\end{proof}

The monoids of type $\CP^2$ constructed in \cite{BB} have $1795$ and $122$
elements. The run that found $M_{23}$ also reports that no monoid with fewer
than $23$ elements has a classifying space of this type; that statement
rests on exhaustive computations which are not reproduced here, and it is
not asserted in this paper.

\section{Provenance and checks}\label{sec:provenance}

This paper is AI-produced. No human has proved, written or checked any part
of it.

\paragraph{What is whose.} Theorem~\ref{thm:BB}, Conjecture~7.1, the map
$j_e$ through the Karoubi envelope and the surjectivity of $\eta$ are
Ben-Bassat's \cite{BB}; Theorem~B of \cite{Mi} is Miller's; (K3) is
Atiyah's. The monoid $M_5$, the statement on its classifying space and the
suggestion that became Theorem~\ref{thm:BB} are Fiedorowicz's \cite{Fie}, and
so is the first published proof of that statement \cite{Fie02};
Lemma~\ref{lem:known}\,(a) is McDuff's \cite{McD}. The statement that a
finite monoid whose classifying space is
homotopy equivalent to a finite-dimensional complex has trivial universal
group, with its proof by Lemmas~\ref{lem:known} and~\ref{lem:je} and
Miller's theorem, and the resulting failure of Conjecture~7.1, were found by
a run of GPT-6 Astra (OpenAI, ChatGPT Pro mode) on 10 November 2026 (run
CATHOT, statements G.4 and G.6), in answer to a research prompt on
Grothendieck's \emph{Les D\'erivateurs}. Theorems~\ref{thm:K},
\ref{thm:main} and~\ref{thm:equiv} and
Corollaries~\ref{cor:finitepi}, \ref{cor:swan} and~\ref{cor:class}, and
Corollary~\ref{cor:cd} with the conclusion $U(M)=1$, were proved by Claude
Fable~5.2 (Anthropic) at effort max on 12 November 2026, while checking that
run. The present form of Corollary~\ref{cor:cd} with its derivation from
\cite{GP}, the last sentence of Remark~\ref{rem:miller}\,(a) and the
question in Remark~\ref{rem:cd} come from the search of the literature made
for this paper on 25 November 2026 by a second Claude Fable~5.2 at effort
max; the writer checked the derivation and wrote the proof. The observation
on the map $j_e$ at the end of Remark~\ref{rem:cd} and the remark that the
case $a=g$ of \cite[Example~3.4]{BB} has trivial universal group and
unbounded homology come from the check of 25 November 2026 described below.
Lemma~\ref{lem:known}\,(c) is known. The author of \cite{BB} was informed of
the counterexample on 12 November 2026.
Proposition~\ref{prop:M6}\,(c), with the resolution used for (b), is due to
the run S2COUNT of GPT-6 Astra (ChatGPT Pro mode), completed on 16 November
2026, which treats a family of monoids containing $M_6$; (b) is Sweeney's.
The monoid $M_{23}$, its resolution, the cocycle $u$ and
Proposition~\ref{prop:M23} are due to the run FINMON of GPT-6 Astra
(ChatGPT Pro mode), 12--17 November 2026. That run was given
Theorems~\ref{thm:main} and~\ref{thm:equiv} as statements to prove or
refute, with the route through $K$-theory suggested, and reports proofs of
them and the observations of Remarks~\ref{rem:miller}\,(b)
and~\ref{rem:rational}; it also treats finite categories with several
objects, which are outside this paper.

\paragraph{The text.} The text was written by Claude Fable~5.2 at effort
max on 25 November 2026. The proof of Theorem~\ref{thm:K} is given here in a
form that uses only finite skeleta and that is local at one prime
throughout; the proofs of Section~\ref{sec:examples} follow those of the two
runs, shortened.

\paragraph{Checks.} The statement of the run CATHOT and its proof were
checked in full by Claude Fable~5.2 at effort max on 12 November 2026,
against the texts of \cite{BB}, \cite{Mi} (with its correction) and
\cite{At}. Theorems~\ref{thm:main} and~\ref{thm:equiv}, in the form of 12
November, were derived again by a second Claude Fable~5.2 at effort max and
were the object of a targeted check by Claude Opus~5.5 at effort high; the
run FINMON, of another model family, reports proofs along the suggested
route, which the writer has not read. The present form of the proof of
Theorem~\ref{thm:K}, Lemma~\ref{lem:je}, the proofs of
Theorems~\ref{thm:main} and~\ref{thm:equiv} and of Corollary~\ref{cor:cd},
and the two examples were checked on 25 November 2026 by Claude Opus~5.5 at
effort max, called the checker below, which recomputed the examples with a
program of its own, independent of the writer's program and of the programs
of the runs. The checker found no error in the proofs; at its indication
the last clause of Corollary~\ref{cor:conj} was corrected and several
citations were made more precise. The checker is a model of the same family
as the writer. A model of another family, GPT-6 Astra (OpenAI, ChatGPT Pro
mode), checked the results in the run FINMON\_CP2 of 25 November 2026 (its
checkpoint of 06:12~UTC). It was given the statements and proofs of 12
November and recorded which of its derivations came before its reading of
them. Before reading them it derived, by routes of its own,
Theorem~\ref{thm:main} for integral homology, Theorem~\ref{thm:equiv},
Corollary~\ref{cor:finitepi}\,(b), the group $eMe$ with the map $j_e$ of
Lemma~\ref{lem:je}, and Corollary~\ref{cor:class}\,(b) from
Theorem~\ref{thm:BB}: the rank-zero classes in $K$-theory are nilpotent by
an argument with the Postnikov sections of $BU$ which needs no finiteness of
the complex, and the finite model in Theorem~\ref{thm:equiv} comes from
\cite[Proposition~4C.1]{Ha}. It then checked the statements for $\F_p$ and
$\Zp$ in Theorem~\ref{thm:main} and Corollary~\ref{cor:finitepi}\,(a)
against the proofs it was given, and Corollary~\ref{cor:conj}, and found no
fault. It has not read the present text. The
paragraph on the monoid $M_5$ in Section~\ref{sec:examples} and the
construction of the classifying space in Remark~\ref{rem:miller}\,(b) were
added after that check; the checker checked both the same day, recomputing
the statements on $M_5$ with its own program, and a second program of the
writer confirms them as well. For
Section~\ref{sec:examples} the writer ran a program of its own,
independent of the programs of the runs, in exact integer arithmetic: for
$M_6$ it confirms the resolution in degrees $0$ to $3$ and, for
$2\le n\le5$, that $z_n$ is a cycle, that $f_n$ is a cocycle and that
$f_n(z_n)=1$; for $M_{23}$ it confirms every finite verification named in
the proof of Proposition~\ref{prop:M23}, with the values $-1$ and $1$ of the
two evaluations. The proof of Theorem~C of \cite{BB}, an unrefereed
preprint, was read in full by the writer on 12 November 2026, without
finding a gap. Nothing is formalized.

\paragraph{Citations.} The writer has opened \cite{BB} (source of the first
version), \cite{Mi} (pp.~39--41 and 84--85, and the correction),
\cite{At} (\S\S6--8), \cite{Sw} (\S\S1, 2, 5), \cite{PW} (introduction),
\cite[Theorem~6.2]{CC} and
the statement of \cite[Proposition~4C.1]{Ha}. The reader who made the search of the
literature has opened \cite{Swan} in full and \cite{St}, \cite{MS},
\cite{GP} and \cite{GS} at the places cited. The checker opened every
reference at the places cited, among them the places of \cite[\S2]{At},
\cite{Ha} and \cite{Qu} that the writer had cited from memory, with the
exception of \cite{AH}, \cite{Fie} and \cite{McD}; the locator in
\cite{AH} is still given from memory. The session that coordinated this
work, a second Claude Opus~5.5 at effort max, opened \cite{Fie} at
pp.~321--324 and \cite{McD} in full on 25 November 2026: the statements
cited from them are as given here, and Theorems~\ref{thm:main},
\ref{thm:equiv} and~\ref{thm:K} are neither on these pages of \cite{Fie} nor
in \cite{McD}. The search for
earlier occurrences of these theorems, with its queries and with what could
not be opened, is part of the records. The proof of Corollary~\ref{cor:cd}
has been read by the writer, by the reader who proposed it and by the
checker.

\paragraph{Records.} This paper is accompanied in its workbench by its
records:
{\raggedright
\begin{itemize}
\item \texttt{CHECK\_EVENTS.tsv}: every check, with model, effort and date;
\item \texttt{PROVENANCE.md}: every contributor and the state of every
citation;
\item \texttt{CHECK\_D2\_INDEPENDENT.md}: the report of the check of 25
November 2026;
\item \texttt{d2\_examples\_check.py} and \texttt{d2\_m5\_check.py}, in the
folder \texttt{checks}: the programs of the writer's checks of
Section~\ref{sec:examples}, with their outputs.
\end{itemize}
\noindent The workbenches are reached from
\url{https://github.com/KokunoYumeto/mathematics-workbenches}.\par}

\begin{thebibliography}{99}
\bibitem{At} M.~F.~Atiyah, Characters and cohomology of finite groups,
 Publ.\ Math.\ Inst.\ Hautes \'Etudes Sci.\ 9 (1961), 23--64.
\bibitem{AH} M.~F.~Atiyah and F.~Hirzebruch, Vector bundles and homogeneous
 spaces, in \emph{Differential Geometry}, Proc.\ Sympos.\ Pure Math.\ 3,
 Amer.\ Math.\ Soc., 1961, 7--38.
\bibitem{BB} O.~Ben-Bassat, On a conjecture of Fiedorowicz,
 arXiv:2610.38760v1, 28 October 2026.
\bibitem{CC} I.~Carcac\'{\i}a-Campos, How many points does it take to see the
 homotopy of a space? Minimal acyclic categorical models, arXiv:2611.09718v1,
 6 November 2026.
\bibitem{Fie} Z.~Fiedorowicz, Classifying spaces of topological monoids and
 categories, Amer.\ J.\ Math.\ 106 (1984), no.~2, 301--350.
\bibitem{Fie02} Z.~Fiedorowicz, A counterexample to a group completion
 conjecture of J.~C.~Moore, Algebr.\ Geom.\ Topol.\ 2 (2002), 33--35;
 arXiv:math/0202260.
\bibitem{GP} R.~Gray and S.~J.~Pride, Homological finiteness properties of
 monoids, their ideals and maximal subgroups, J.\ Pure Appl.\ Algebra 215
 (2011), no.~12, 3005--3024; arXiv:1003.3227v1.
\bibitem{GS} R.~D.~Gray and B.~Steinberg, Topological finiteness properties
 of monoids. Part~1: Foundations, Algebr.\ Geom.\ Topol.\ 22 (2022),
 no.~7, 3083--3170; arXiv:1706.04387v1, whose numbering is used here.
\bibitem{Ha} A.~Hatcher, \emph{Algebraic topology}, Cambridge University
 Press, 2002.
\bibitem{MS} A.~Martinian and B.~Steinberg, Realizing wedges of Moore spaces
 as classifying spaces of finite semigroups, Semigroup Forum (2026),
 \url{https://doi.org/10.1007/s00233-026-10648-0}; arXiv:2511.12657v1, whose
 numbering is used here.
\bibitem{McD} D.~McDuff, On the classifying spaces of discrete monoids,
 Topology 18 (1979), 313--320.
\bibitem{Mi} H.~Miller, The Sullivan conjecture on maps from classifying
 spaces, Ann.\ of Math.\ (2) 120 (1984), 39--87; Correction, Ann.\ of Math.\
 (2) 121 (1985), 605--609.
\bibitem{PW} M.~Palmer and X.~Wu, Embedding groups into acyclic groups,
 arXiv:2510.16879v1, 19 October 2025.
\bibitem{Qu} D.~Quillen, Higher algebraic $K$-theory~I, in \emph{Algebraic
 $K$-theory~I}, Lecture Notes in Math.\ 341, Springer, 1973, 85--147.
\bibitem{St} B.~Steinberg, The homology of completely simple semigroups,
 Semigroup Forum 110 (2025), 745--761; arXiv:2405.06594v1, whose numbering is
 used here.
\bibitem{Swan} R.~G.~Swan, The nontriviality of the restriction map in the
 cohomology of groups, Proc.\ Amer.\ Math.\ Soc.\ 11 (1960), 885--887.
\bibitem{Sw} D.~Sweeney, Homotopy types of small semigroups,
 arXiv:2502.05276v2, 13 June 2026.
\end{thebibliography}
\end{document}
